\cotefile{157-5}
\scanpage{277}
\mspage{193}
% archive sheet number: 276
\diagram{\begin{tikzcd}[ampersand replacement=\&,column sep=small]
C \arrow[rr,"\alpha"] \arrow[dr,dash,"i"'] \& \& C' \arrow[dl,"i'"]\\
\& E \&
\end{tikzcd}}

% slanted
\margin{$*$ En fait, il suffit qu'il satisfasse la condition plus faible a) du $1^\circ)$.}

le foncteur $i$ est strictement générateur\textsuperscript{$*$}, et si de plus $i'$ est pleinement fidèle, alors $i'$ est strictement générateur.

% slanted
\margin{\unsure{permuter|renv. de|ren. du} \unsure{n. et|n. 2 et} $3^\circ$ \unsure{$\big($}}

\deleted{\unsure{3|2}}$\unsure{2|3}^\circ)$ Supposons que $i$ soit strictement générateur.\textsuperscript{$*$} Alors \inserted{pour \unsure{tt} objet $X$ de $M$,} la famille des flèches $i(x)\longrightarrow X$, indexée \unsure{sur|par} $\mathrm{Ob}\,C_{/X}$, est non seulement épimorphique stricte (condition a) du $1^\circ)$, mais \inserted{de plus} \deleted{\unsure{t|a}}\unsure{toute|tout} \unsure{famille|familles} de flèches de but $X$, incluant la famille précédente, est également épimorphique stricte.

% slanted
\margin{$*$ Il suffit $=$ qu'elle satisfasse la condition plus faible a) du $1^\circ)$.}

\deleted{\uline{\unsure{Démonstration de}} 4.8.2.}

\uline{Corollaire} 4.8.3. Soit $M$ une $\unsure{\mathcal U|U}$-catégorie, $C$ une sous-catégorie strictement génératrice\textsuperscript{$*$}, $C'$ une sous-catégorie \uline{pleine} de $M$ qui contient $C$, alors $C'$ est stricte\deleted{\unsure{ment}} génératrice.

% slanted
\margin{\deleted{\unsure{il suffit}}}

C'est en effet un cas particulier de 4.8.2, $3^\circ$\unsure{.}

\scanpage{278}
\mspage{194}
% archive sheet number: 277
\uline{Dém. de} 4.8.2. La démonstration est \inserted{du $1^\circ)$ est} celle de SGA 4 I 7.2. (i). (\unsure{Où|où} on suppose que $C$ est une sous-catégorie pleine de \deleted{$\unsure{C}$}$M$, mais la démonstration donnée prouve en fait l'énoncé donné ici. ) Prouvons $2^\circ)$. Rajoutons donc à la famille des

\[ i(x)\longrightarrow X \qquad (x\in\mathrm{Ob}\,C_{/X} \]

une famille

\[ X_{\alpha}\xrightarrow{p_{\alpha}}X \qquad (\alpha\in A) \]

de nouvelles flèches, et prouvons que la famille totale ainsi obtenue est \unsure{une|encore} épimorphique stricte. On doit donc prouver que si $Y$ est un objet de $M$, et si on a une famille de flèches $u_x:i(x)\longrightarrow Y$, $u_\alpha:Y_\alpha\longrightarrow Y$, satisfaisant la condition \deleted{$\unsure{Y_\alpha|X_\alpha}$} de compatibilité habituelle pour des \inserted{couples de} flèches

\diagram{\begin{tikzcd}[ampersand replacement=\&,column sep=small,row sep=small]
\& \& \deleted{\unsure{x|i}}i(x) \arrow[dr,dashed,"p_x"] \& \\
(v_x,v_y): \& Z \arrow[ur,"v_x"] \arrow[dr,"v_y"'] \& \& X\\
\& \& i(y) \arrow[ur,dashed,"p_y"'] \&
\end{tikzcd}}

ou

\diagram{\begin{tikzcd}[ampersand replacement=\&,column sep=small,row sep=small]
\& \& i(x) \arrow[dr,dashed,dash,"p_x"] \& \\
(v_x,\unsure{v_\alpha|v_a}) \& Z \arrow[ur,"v_x"] \arrow[dr,"v_\alpha"'] \& \& \unsure{Y|X}{,}\\
\& \& X_\alpha \arrow[ur,dashed,dash,"p_\alpha"'] \&
\end{tikzcd}}

\scanpage{279}
\mspage{\unsure{195}}
% archive sheet number: 278
ou

\diagram{\begin{tikzcd}[ampersand replacement=\&,column sep=small,row sep=small]
\& \& \deleted{\unsure{Y}}X_\alpha \arrow[dr,dashed,dash] \& \\
(v_\alpha,v_\beta) \& Z \arrow[ur,"v_\alpha"] \arrow[dr,"v_\beta"'] \& \& X{;}\\
\& \& X_\beta \arrow[ur,dashed,dash] \&
\end{tikzcd}}

savoir les relations

{\makeatletter\tagsleft@true\makeatother
\[
\deleted{u_xv_x=u_yv_y}\quad
\left\{\begin{aligned}
&\unsure{\mathrm a)}\ \deleted{\unsure{p_x}}u_xv_x=u_yv_y\\
&\unsure{\mathrm b)}\ u_xv_x=u_\alpha v_\alpha\\
&\unsure{\mathrm c)}\ u_\alpha v_\alpha=u_\alpha v_\beta,
\end{aligned}\right. \tag*{$(\alpha)$}
\]
}

alors il existe une flèche \deleted{\unsure{$u:X$}}

\[ u:X\longrightarrow Y \]

(\unsure{néc.|nèc.} unique, d'après la propriété a) d'épimorphicité stricte de la famille initiale indexée par $C_{/\unsure{X|x}}$) telle qu'on ait

{\makeatletter\tagsleft@true\makeatother
\[
\left\{\begin{aligned}
up_x&=v_x &&\text{pour }\unsure{\mathrm{tt}|\forall}\ x\ \unsure{\in}\mathrm{Ob}\,C_{/X}\\
up_\alpha&=v_\alpha &&\text{pour }\mathrm{tt}\,\alpha\in A.
\end{aligned}\right. \tag*{$(*)$}
\]
}

Or d'après la propriété a), il existe une unique $u:X\longrightarrow Y$, de telle façon que la première des relations (*) soit satisfaite, pour $\mathrm{tt}\,x\in\mathrm{Ob}\,C_{/X}$. Donc il s'agit seulement de prouver que pour cet $u$, on a \deleted{\unsure{pour tout}} $\alpha$ la relation

\scanpage{280}
\mspage{196}
% archive sheet number: 279
\[ up_{\alpha}=u_{\alpha}. \]

Or, en appliquant la propriété a) à $X_\alpha$ au lieu de $X$, on voit que la famille des flèches

\[ i(\deleted{\unsure{x}}y)\xrightarrow{\deleted{\unsure{p_y}}q_y}X_\alpha\ \unsure{,|\phi} \]

(indexée par $C_{/X_\alpha}$) est épim. stricte, a fortiori épimorphique. Donc pour prouver $up_\alpha=u_\alpha$, il suffit de prouver qu'on a

{\makeatletter\tagsleft@true\makeatother
\[ (up_\alpha)q_y=u_\alpha q_y \tag*{$(**)$} \]
}

pour \unsure{tt|tout} $y$\deleted{\unsure{$\in\mathrm{Ob}$}} dans $C_{/X_\alpha}$. Mais c'est là une conséquence de la compatibilité ($\alpha$) (\unsure{deuxième|dernière} \unsure{formule}), appliquée au cas suivant\unsure{:}

\diagram{\begin{tikzcd}[ampersand replacement=\&,row sep=small]
\& i(y) \arrow[dr,dashed,"p_\alpha q_y=p_x"] \&\\
Z=i(y) \arrow[ur,"\mathrm{id}"] \arrow[dr,"q_y"'] \& \& X{,}\\
\& X_\alpha \arrow[ur,dashed,"p_\alpha"'] \&
\end{tikzcd}}

on a donc, en posant $\deleted{\unsure{x_\alpha}}x=(y,p_\alpha q_y)\in\mathrm{Ob}\,C_{/X}$,

\[ \underbrace{u_{\unsure{y|x}}\circ\mathrm{id}_{i(y)}}_{u_x}=u_\alpha q_y, \]

donc pour prouver $\unsure{(**)|(\alpha)}$, il reste à vérifier

\[ u_x=\underbrace{(u\deleted{\unsure{q}}p_\alpha)q_y)}_{u(p_\alpha q_y)=up_x}\deleted{\unsure{\mathrm{id}}}, \]

\scanpage{281}
\mspage{197}
% archive sheet number: 280
qui est en effet \unsure{vérifiée|vérifié} par définition de $u$. Cela \unsure{achève|qchève} de prouver $2^\circ$.

Prouvons $3^\circ$. Comme $i'$ est pl. fidèle, en vertu du $1^\circ$, pour prouver que $i'$ est str. générateur, il suffit d'établir que \inserted{pour \unsure{tt} objet $X$ de $M$,} la famille des flèches $i'(x')\longrightarrow X$, indexée par $\mathrm{Ob}\,C'_{/X}$, est épimorphique stricte.

Or cette famille contient la famille des $\underbrace{i'(\alpha(x))}_{\substack{\Vert\\i(x)}}\longrightarrow X$, \deleted{\unsure{$x$}} \inserted{indexée par $\mathrm{Ob}\,C_{/X}$.} \deleted{\unsure{cette dernière étant épi}} Donc il résulte de l'hypothèse sur $i$, et du $2^\circ$, que cette famille est épimorphique stricte. On a gagné !

Ouf\unsure{...|.--} \par\hfill\unsure{$\backslash$}\hfill\null

\uline{Démonstration de 4.8.1}. On a

\[ c'\Longrightarrow c\Longrightarrow b\Longrightarrow a, \]

la première implication \unsure{tautologique}, la seconde résulte du fait que $\unsure{\pi|\kappa|k}_0$ \unsure{cardinal|card.} régulier pour la fega canonique, implique que

\[ M\xleftarrow{\approx}\mathrm{Ind}_{\pi_0}(\unsure{F|\mathcal F}^{\pi_0}), \]

$F^{\pi_0}$ équivalente à une petite catégorie \inserted{$C$}, d'où b). \deleted{\unsure{ou}}

\scanpage{282}
\mspage{198}
% archive sheet number: 281
Enfin b) $\Longrightarrow$ a) résulte du fait que $\mathrm{Ind}_{\unsure{\pi|\kappa|k}_0}(C)$ est stable par $\varinjlim$ \unsure{grdes|gdes} devant $\pi_0$ (4.4.15), et que $C$ est une sous-catégorie str. génératrice de $\mathrm{Ind}_{\pi_0}(C)$, formée d'objets $\pi_0$-accessibles.

\deleted{Pour vérifier que $C$ est strictement génératrice, en vertu de 4.8.2, $1^\circ$, il suffit, puisque $C$ est une sous-catégorie \uline{pleine} de $M=\mathrm{Ind}_{\pi_0}(C)$, de vérifier que pour \unsure{tt} \inserted{$F$ dans $M$, la} famille des flèches \deleted{de} $x\longrightarrow F$ de but $F$, de source dans $C$, est épim. stricte. Pour ceci, \deleted{il suffit d'ailleurs de le vérifier \unsure{dans}}}

\unsure{Il|il} nous faut expliciter le

\uline{Lemme} 4.8.4. Soit $C$ une \inserted{$\unsure{\mathcal U|U}$-}catégorie, \deleted{\unsure{égale|équi.}} \inserted{$\pi_0$ un cardinal infini} \deleted{\unsure{à une} petite} Alors $C$ est \unsure{strictement|stricte} génératrice dans \inserted{$M=$}$\mathrm{Ind}\, C$ et dans \inserted{$M=$}$\mathrm{Ind}_{\pi_0}C$.

\deleted{Il faut prouver, $F$ étant un objet de $M$, que}

\[ \deleted{\varinjlim_{C_{/F}}x\longrightarrow F} \]

\deleted{est un iso. On a $F=\varinjlim_{\deleted{\unsure{x}}{\unsure{I|\underline I}}}x_i$, $I$ \inserted{ens. ordonné} filtrant (\unsure{grd} devant $\pi_0$ dans le deuxième cas), les $x_i$ dans $C$, pour un foncteur $I\longrightarrow C\deleted{\unsure{_{/F}}}$ convenable. On a une factorisation \inserted{can.} \unsure{de}}

\scanpage{283}
\mspage{199}
% archive sheet number: 282
\deleted{ce foncteur par $C_{/\unsure{x|X}}$,}

\deleted{\diagram{\begin{tikzcd}[ampersand replacement=\&,column sep=large]
I\arrow[r] \& C_{/X}\arrow[r,dashed] \& C\arrow[r,dashed] \& M{,}
\end{tikzcd}}}

\deleted{et on a vu (\unsure{cf} 4.4.2 p.~\unsure{74|79|71}) que ce foncteur est cofinal.}

En effet, le foncteur canonique $M\longrightarrow C^{\wedge}$ est pleinement fidèle, et on applique le critère \unsure{c)} de 4.8.2.

Pour achever de prouver 4.8.1, il reste donc à établir l'implication

\[ \mathrm a)\Longrightarrow c'. \]

\deleted{Soit donc $S\subset M$}

Supposons donc que $M$ satisfasse à $L_{\unsure{\pi|\kappa|k}_0}$, et admette une petite sous-catégorie str. génératrice $S$, formée d'objets $\pi_0$-accessibles. \deleted{L'hyp. $L_{\pi_0}$ \unsure{signifie|dit} que} \inserted{On} \unsure{trouve|tient} un foncteur canonique

\[ \mathrm{Ind}_{\pi_0}S\longrightarrow M, \]

prolongeant l'inclusion \deleted{canonique} $S\hookrightarrow M$, de façon à commuter aux $\varinjlim$ \unsure{grdes} devant $\pi_0$, et du fait que $S\subset M_{\pi_0}$, ce foncteur est \uline{pleinement fidèle}. Mais comment prouver qu'il est \unsure{ess. surjectif} ? \unsure{Ca|Ce} ne sera d'ailleurs pas vrai \deleted{\unsure{à bien}} sans autre hypothèse sur $S$. \deleted{\unsure{en fait}} L'idée naturelle est de saturer $S$ pour les $\varinjlim$

\scanpage{284}
\mspage{200}
% archive sheet number: 283
de cardinal $\leq\unsure{\pi|\kappa|k}_0$, représentables dans $M$, et notamment par fact. \unsure{directs}, \deleted{\unsure{grdes} devant $\pi_0$}, mais même alors, et $=$ en admettant que $M_{\pi_0}$ lui-$=$\unsure{!} est équivalent à une petite catégorie\deleted{\unsure{)}} , et prenant $C=M_{\pi_0}$\deleted{\unsure{)}} , il n'est pourtant pas clair que \deleted{\unsure{le f.}} \inserted{ce fonct. pl. fid.} $\mathrm{Ind}_{\pi_0} C\longrightarrow M$ soit ess. surjectif), ce qui équiv\deleted{\unsure{alvait|alait}} revient à dire que pour \unsure{tt|$\forall$} $\deleted{\unsure{E}}F$ dans $M$, $C_{/F}$ est \unsure{grde} devant $\pi_0$. Il n'est en effet pas $=$ clair que $C_{/\unsure{F|P}}$ soit filtrante.

Donc finalement, je doute qu'on ait a) $\Longrightarrow$ c') ou seulement a) $\Longrightarrow$ b), la situation n'est pas bien comprise visiblement.

Mais prouvons quand même

\[ b\Longrightarrow c'. \]

Pour ceci, \unsure{OPS|WPS} \deleted{\unsure{par}} $M=\mathrm{Ind}_{\pi_0}C$, $C$ petite. On sait alors que la filtration de $M$ par les $(M_{\unsure{\pi_0|\pi}})_{\pi\geq\pi_0}$ est une \deleted{\unsure{filtrati}} fega (prop. 4.7.7 et déf. \unsure{4.7.8'|4.7.8}), et $=$ une feg \uline{canonique}. De plus, je \inserted{dis que} le cardinal $\pi_0$ est \uline{utile} \inserted{et $=$ régulier} pour cette filtration. Ce qui achèvera de prouver c'). \deleted{En effet,} On a un résultat un peu plus précis\unsure{:}

\scanpage{285}
\mspage{201}
% archive sheet number: 284
\uline{Corollaire} 4.8.5. Soit $C$ une \deleted{petite} catégorie équivalente à une petite catégorie, \deleted{$\unsure{M}$} $\unsure{\pi|\kappa|k}_0$ un cardinal infini, $M=\mathrm{Ind}_{\pi_0}(C)\subset C^{\wedge}$. Alors

$1^\circ)$ $\pi_0$ est un cardinal \deleted{\uline{régulier}}\inserted{utile, et $=$ \uline{régulier}} pour la \deleted{\unsure{filtr}} feg canonique $(M_\pi)_{\pi\geq\pi_0}$ de $M$.

$2^\circ)$ $M_{\pi_0}\unsure{=|\simeq}C^{\unsure{\natural|4}}$ (sous-catégorie pleine de \inserted{$M=$}$\mathrm{Ind}_{\pi_0}C$ formée des facteurs directs dans $M$ des objets de $C$), et $M_{\pi_0}$ est une enveloppe de \unsure{Karoubi} de $C$. \unsure{a}

% slanted
\margin{\uline{4.8.5.2} \unsure{Expliciter|Explicitons} aussi ceci : Soit $u:C\longrightarrow C'$ un foncteur, \unsure{avec|où} $C,C'$ \unsure{deux|sont} petites catégories. Conditions équivalentes\unsure{:} \par a) $u$ induit $\mathrm{Kar}\,C\xrightarrow{\approx}\mathrm{Kar}\,C'$. \par a') $u$ pl. fid., et tout objet de $C'$ \unsure{isom.} à un facteur \unsure{direct|dir.} d'un \unsure{objet} de l'image essentielle\unsure{.} \par b) $\mathrm{Ind}(u):\mathrm{Ind}(C)\longrightarrow\mathrm{Ind}(C')$ est une équiv. \par c) $\mathrm{Ind}_{\pi_0}(u):\mathrm{Ind}_{\pi_0}(C)\longrightarrow\mathrm{Ind}_{\pi_0}(C')$ est une équivalence ($\pi_0$ un cardinal infini donné) \par d) $u_!:C^{\wedge}\longrightarrow C'^{\wedge}$ est une équivalence \par d') $u_*:C'^{\wedge}\longrightarrow C^{\wedge}$ est une équivalence \par d'') $u^{\unsure{\wedge|*}}:C'^{\wedge}\longrightarrow C^{\wedge}$ est une équivalence}

\margin{\uline{Corollaire} 4.8.5.1. Si $C$ est Karoubienne, alors $M_{\pi_0}\subset C^{\wedge}$ est la catégorie formée des foncteurs \uline{représentables}, donc $C\xrightarrow{\approx}M_{\pi_0}$.}

\uline{Dém} \inserted{$1^\circ)$} On sait déjà que

\[ \mathrm{Ind}_{\pi_0}M_{\pi_0}\longrightarrow M \]

est pleinement fidèle, et pour voir que c'est ess. surjectif, on note que l'on a

\[ M=\mathrm{Ind}_{\pi_0}C\mathrel{\unsure{\leftrightarrow|\hookrightarrow}}\mathrm{Ind}_{\pi_0}M_{\unsure{0|\pi_0}}\longrightarrow M, \]

le composé étant isom. \deleted{\unsure{donc}} au foncteur identique, est ess. surjectif, donc aussi \unsure{cas.|ceci.} Cela prouve donc que $\pi_0$ est un cardinal utile pour la filtration envisagée, prouvons que ce cardinal est $=$ régulier, donc que pour

\[ F\in\mathrm{Ob}\,M_{\pi},\qquad\text{avec }\left[\begin{aligned}
\pi\ \deleted{\unsure{\geq}}&>\pi_0\\
\pi^{\pi_0}&\geq\pi_1,
\end{aligned}\right. \]

on a

\[ F=\varinjlim_I \deleted{\unsure{X}}x_i \quad \left|\begin{aligned}
&x_i\in\mathrm{Ob}\,M_{\pi_0}\\
&I\ \text{grd devant }\pi_0\\
&\mathrm{Card}\, I\leq\pi^{\pi_0}.
\end{aligned}\right. \]

\deleted{\unsure{pour} $\pi^{\pi_0}\geq\pi_1$\unsure{.}}

\scanpage{286}
\mspage{202}
% archive sheet number: 285
Pour ceci, on utilise (4.7.6.2) page 171

\[ M_{\unsure{\pi|\kappa|k}}\subset\mathrm{Filt}^{\pi^{\pi_0}}\qquad(\text{si }\pi\geq\pi_0,\pi^{\pi_0}\geq\pi_1), \]

et la définition de $\mathrm{Filt}^{\pi'}$ (pour $\pi'\geq\pi_0$) \deleted{\unsure{donne}} montre que c'est là la condition cherchée.

Prouvons donc $2^\circ)$. Comme $M_{\pi_0}$ est stable par \inserted{les} $\varinjlim$ inductives finies qui sont représentables dans $M$ (SGA 4 I \unsure{9.9.}), il est stable par facteurs directs, et comme il contient $C$, on a donc $C^{\unsure{\natural|4}}\subset M_{\pi_0}$. \deleted{\unsure{Pour l'i}} L'inclusion inverse $M_{\pi_0}\subset C^{\natural}$ se voit de la façon \unsure{standard|standart}, en écrivant un objet \inserted{$F$} de $M_{\pi_0}$ comme $\varinjlim_I x_i$, les $x_i$ dans $C$, $I$ grd devant $\pi_0$, alors $F$ \deleted{\unsure{se}} est facteur direct d'un des \unsure{$x_i$|ci}.

Le fait que $C^{\natural}$ est une \uline{enveloppe de Karoubi} de $C$ dans $M\hookrightarrow C^{\wedge}$ provient du fait que $M$ est stable par facteurs directs, et que $C$ est une sous-catégorie pleine de $M$, cf SGA 4 I (?...).

4.8.6. \inserted{\uuline{\unsure{Questions ouvertes|Question ouverte}}} On a ainsi prouvé 4.8.1 revu et corrigé. Il reste les questions \deleted{de}

$1)$ \unsure{Donner} un exemple où \inserted{4.8.1} a) est \unsure{bien|elle} vérifiée, mais non b) etc.

% slanted
\margin{Il suffit de trouver un cardinal non régulier \deleted{\unsure{pour}} pour \unsure{la} feg \unsure{canon.} de $\mathrm{Ind}_{\pi_0}C$.}

$2)$ Dégager des conditions supplémentaires \inserted{\unsure{serviables|vérifiables|raisonnables}} qui, jointes à a), impliquent b), donc c, c'.

\scanpage{287}
\mspage{203}
% archive sheet number: 286
Signalons aussi pour mémoire les questions suivantes :

3) Soit $M=\mathrm{Ind}_{\unsure{\pi|\kappa|k}_0}(C)$, et $\pi>\pi_0$. Alors $\mathrm{Filt}^{\pi}$ est-il stable par facteurs directs ? Cela est-il vrai du moins si $\pi=\pi^{\pi_0}$, ou ce qui revient au $=$ (cf. 4.7.3), a-t-on dans ce cas $\mathrm{Filt}^{\pi}=\mathrm{Filt}'^{\pi}$ ? Si oui, il serait inutile dans la description 4.7.7. des filtrations \unsure{cardinales|canoniques|coïncidentes}, avec la majoration de cardinaux (4.7.7.1) $\mathrm{Card}\,I\leq\pi'^{\pi}$, de mettre l'hypothèse restrictive $\pi'^{\pi}\geq\pi_1$.

Pour avoir une réponse affirmative à cette question, pour $\pi$ donné, il suffi\deleted{\unsure{ait}}t \inserted{(et probablement faudrait)} d'avoir une réponse affirmative à celle-ci : soient $D,D_0$ deux catégories \inserted{\deleted{\unsure{filtrantes}} \unsure{grdes} devant $\pi_0$}, $D_0$ facteur direct de $D$, $D$ admettant une catégorie cofinale $I$ grde devant $\pi_0$, de cardinal $\leq\pi$. Alors en est-il de $=$ de $D_0$ ? \deleted{Si $p$ est une rétraction} de $D$ sur $D_0$, la difficulté est que je ne sais pas si $p$ est \inserted{toujours} cofinal, \deleted{\unsure{Mais en}} fait, je \unsure{vois|crois} que

\deleted{la réponse est négative, Soit en effet $D_0$ une cat. grde devant $\pi_0$, et n'admettant \unsure{pas} un \deleted{\unsure{foncteur}} cofinal $I_0\longrightarrow D_0$, avec \inserted{$I_0$ filtrant et $=$} $\mathrm{Card}\,I_0\leq\pi$ (\unsure{a fortiori|le foncteur}, avec $\mathrm{Card}\,I_0\leq\pi$ et $I_0$ grde devant $\pi_0$). Par exemple, on peut prendre l'ens. ordonné $\unsure{I|\underline{I}|\overline{1}}_\alpha$ des ordinaux $<\alpha$, où $\alpha$ est le plus petit ordinal tel que $\mathrm{Card}\,\alpha>\pi$. Soit $D$ la catégorie déduite de $D_0$ en ajoutant un élément final (\uline{cône} droit sur $D_0$). Alors $D$ admet \deleted{une} \inserted{un} \unsure{cofoncteur|foncteur} cofinal $e\longrightarrow D$, où $e$ est grd devant $\pi_0$ et de cardinal $\leq\unsure{\pi_1|\pi}$, et}

\scanpage{288}
\mspage{204}
% archive sheet number: 287
\deleted{$D_0$ est un facteur direct de $D$},

ce n'est pas tjrs le cas, $=$ si \deleted{$\unsure{D}$}$D_0$ est finie. Prenons p. ex. pour $D_0$ une catégorie ayant objet final \deleted{\unsure{et}} $e$, et une flèche $x\longrightarrow e$ qui n'est pas iso, et admettant une section (p. ex. $D_0$ la ss-catégorie pleine de $\mathbb{\Delta}$ formée par $\mathbb{\Delta}^0,\mathbb{\Delta}^1$ :

\diagram{\begin{tikzcd}[ampersand replacement=\&]
\Delta^0 \arrow[r,shift left=1.3ex,"{\unsure{\partial_0|d_0}}"] \arrow[r,"{\unsure{\partial_1|d_1}}" description] \& \Delta^1 \arrow[l,bend left=28,"\unsure{k|r|\sigma}"']
\end{tikzcd}}

). Soit $D$ déduit de $D_0$ en ajoutant un objet final $\widetilde e$. Alors $D_0,D$ sont filtrantes, et on obtient une rétraction de $D$ sur $D_0$, envoyant $\widetilde e$ sur $x$. Ce foncteur n'est pas cofinal, car \deleted{s'il l'était,} un foncteur cofinal doit envoyer $\widetilde e$ sur $e$.

Plus généralement, si $D_0$ est une catégorie telle qu'il existe une flèche

{\makeatletter\tagsleft@true\makeatother
\[ \alpha:e_{D_0^{\wedge}}\longrightarrow x_0 \tag*{$(*)$} \]
}

dans $D_0^{\wedge}$, avec $x_0\in\mathrm{Ob}\,D_0$, alors on voit que $D_0$ est grde devant $\unsure{\pi|\kappa|k}_0$ quel que soit le cardinal infini $\pi_0$, et que d'autre part, désignant encore par $D$\unsure{.} la cat. déduite de $D_0$ en lui ajoutant un objet final \deleted{\unsure{il}} $\deleted{\unsure{x}}e_D$, il y a une rétraction de $D$ sur $D_0$ qui envoie $e_D$ dans $x_0$. Cette rétraction n'est cofinale que si $x$ est un objet final de $\unsure{D|D_0}$, i.e \unsure{si|ds} $\unsure{\alpha}$ est un iso. Supposons qu'il n'en soit pas ainsi.

\deleted{Comme $D$ a un foncteur cofinal $e_D\longrightarrow D$, avec $\mathrm{card}\,(e_D)\deleted{\unsure{\geq}}=1\leq1)$\unsure{$_0$}, on a envie de trouver $\pi>\deleted{\unsure{\pi_0}}\pi_0$, \inserted{$\pi_0=\aleph_0$}, disons \deleted{$\unsure{\aleph_1}$}$\pi=2^{\aleph_0}$, et $D_0$ tels que $D_0$ n'admette pas un foncteur cofinal $I\longrightarrow D_0$, avec $\mathrm{Card}\, I\leq\pi$.}

Mais il y a un foncteur cofinal de la sous-catégorie $I_0$ de $D_0$, telle que $\mathrm{Ob}\,I_0=\{x_0\}$, $\mathrm{Fl}\,I_0=\mathrm{End}_{I_0}(x_0)=\{p,\mathrm{id}_{x_0}\}$, où $p$ est le composé $x_0\xrightarrow{\mathrm{can}}e_{D_0^{\wedge}}\xrightarrow{\alpha}x_0$. Donc pas de contre-exemple de cette façon !

\scanpage{289}
\mspage{205}
% archive sheet number: 288
4) Considérons $M=\mathrm{Ind}_{\unsure{\pi|\kappa|k}_0}(C)$, \deleted{\unsure{$C$ Karoubi}} \inserted{$C$ petite,} et considérons la filtration canonique $(M_\pi)_{\pi\geq\pi_0}$. \deleted{\unsure{Y a-t}} Donner des exemples \deleted{de cardinaux} où il y a des cardinaux non utiles, ou mieux, où pour \unsure{tt|$\forall$} cardinal $\pi$, existe un cardinal $\pi'>\pi$ qui soit non utile. (Cela donnera l'exemple demandé dans $1^\circ$.) D'autre part, donner des exemples de cardinaux utiles qui ne soient réguliers -- si possible, comme tantôt, \deleted{des} tels cardinaux aussi grands qu'on veut. \par\hfill\unsure{$\backslash$}\hfill\null

\uline{Définition} 4.8.\unsure{?|7}. Soit $M$ une $\unsure{\mathcal U|U}$-catégorie, \inserted{$\pi_0$ un cardinal infini.} On dit que $M$ est $\pi_0$-\uline{accessible} si $M$ est équivalente à une catégorie $\mathrm{Ind}_{\pi_0}(C)$, \unsure{-|où} une petite catégorie, i.e si elle satisfait aux conditions équivalentes b, c, c' de 4.8.1 (impliquant la condition a), sans doute strictement plus faible). On dit que $M$ est \uline{accessible} s'il existe un cardinal tel que $M$ \unsure{soit|est} $\pi_0$-accessible [NB si $\pi_0$ est un tel cardinal, alors tout cardinal $\pi>\pi_0$ tel que $\pi^{\pi_0}=\pi$, notamment les cardinaux $\pi'^{\pi_0}$ avec $\pi'\geq2$, ont la $=$ propriété.], ou ce

\scanpage{290}
\mspage{206}
% archive sheet number: 289
qui revient au $=$, \inserted{d'après 4.8.1} si $M$ admet une fega, ou encore une feg canonique.

% slanted
\margin{\uline{\unsure{Commentaires|Commentaire}}}

\deleted{\unsure{4.0}}\unsure{4.8.8}. \deleted{Pour ce qui est de la \unsure{nomenclature}} \inserted{Cette définition est en contradiction avec} \inserted{la terminologie} \deleted{\unsure{suivie}} introduite dans SGA 4 I 9.4 (p 114)\unsure{);|),|;} \deleted{\unsure{q}} \unsure{nous|i.e.} désignons dorénavant sous les noms de ``catégorie \unsure{fbt.|fblt.} $\unsure{\pi|\kappa|k}_0$-accessible'', ``cat. fbt. accessible'' les notions introduites dans 9.4. -- i.e. les ct. admettant une \deleted{\unsure{petite}} \deleted{\unsure{sous-générateurs}} \inserted{sous-catégorie} \unsure{pleines|pleine} \unsure{génératrices|génératrice} (pas $=$ str. génératrice) dont les objets sont $\pi_0$-accessibles (resp. accessibles)

Notons que si $M$ est $\pi_0$-accessible (resp accessible), \inserted{la condition a de} 4.8.1 nous dit que $M$ est \unsure{fb\textsuperscript{=}|fb\textsuperscript{t}|fbt.} $\pi_0$-accessible (resp fbt. accessible), et $=$ qu'il existe une petite sous-catégorie pleine \uuline{strictement} génératrice, qui soit formée d'objets $\pi_0$-accessibles (resp. accessibles). Je doute fort que cette \inserted{dernière} condition \inserted{nécessaire} soit aussi suffisante pour que $M$ soit $\pi_0$-accessible (resp accessible). Mais en vertu de 4.6.3, $3^\circ$ (p 150), il en est ainsi si on suppose $C$ stable par petites $\varinjlim$, ou, ce qui revient au $=$ moyennant $L_{\pi_0}$, \inserted{stable} \unsure{par} limites inductives de cardinal $\leq\pi_0$,

\scanpage{291}
\mspage{\unsure{207}}
% archive sheet number: 290
\uline{Proposition} 4.8.9. Soit $M$ une $\unsure{\mathcal U|U}$-catégorie accessible (4.8.7), \inserted{$\unsure{\pi|\kappa|k}_0$ un cardinal tel que $M$ satisfait à $L_{\pi_0}$.} Conditions équivalentes\unsure{:}

a) $M$ est une catégorie \unsure{protopos|pretopos} (cf 4.6.3, $3^\circ$ p. \unsure{150|152}.)

b) $M$ est stable par petites limites inductives

b') \deleted{(si $M$ satisfait $L_{\pi_0}$)} $M$ est stable par \unsure{plimites|limites} inductives de cardinal $\leq\pi_0$

b'') (Si $M$ est équivalente à une catégorie $\mathrm{Ind}_{\pi_0}(C)$ avec $C$ une \deleted{petite} catégorie \uline{Karoubienne}) \deleted{La catégorie} La catégorie $C$ est stable \inserted{(en elle-même)} par limites inductives de cardinal $\leq\pi_0$.

c) $M$ est stable par petites limites projectives.

% slanted
\margin{\uline{NB} L'\unsure{équivalence|implication} de $c'$, $c''$ \unsure{avec|par} les autres conditions \unsure{n'est} pas établie ici, mais seulement dans 4.11. et plus \unsure{loin}.}

% slanted
\margin{\uline{NB} les conditions $c'$, $c''$ me paraissent plus techniques et moins utiles.}

[ c') (sous les hypothèses $M\approxeq\mathrm{Ind}_{\pi_0}(C)$, et \deleted{Démonstration} quand on se donne un cardinal $\pi$ tel que $\pi\geq\pi^{\pi_0}\geq\pi_1=\mathrm{Sup}(\pi_0,\mathrm{Card}\,\unsure{\mathrm{red}|\text{réd}|\mathrm{rel}}\,C)$ $M_\pi$ \deleted{$\unsure{\leq}$} est stable par $\varprojlim$ finies, et $M$ est stable par petits produits infinis (ou encore, par $\varprojlim$ filtrantes).

c'') \deleted{On a $M\approxeq$} Il existe un cardinal infini $\pi$, et une petite catégorie \deleted{\unsure{D}} stable par $\varprojlim$ finies, telle que $M\approx\mathrm{Ind}_\pi(D)$, de plus $M$ stable par \deleted{produits infinis (ou encore} par $\varprojlim$ filtrantes). ]

\uline{Démonstration}. Les implications suivantes sont tautologiques ou déjà \inserted{$\pm$} connues\unsure{:}

\diagram{\begin{tikzcd}[ampersand replacement=\&,column sep=small,row sep=small]
a \arrow[r,Leftrightarrow,"\unsure{\cdot}"'] \arrow[d,draw=none,"{\deleted{\unsure{\Downarrow|\Updownarrow}}}" description] \& b \arrow[r,Rightarrow,"{\text{\scriptsize cf. 4.4.14}}"] \arrow[dl,Rightarrow] \& b'\unsure{{[}} \arrow[r,Rightarrow,"\unsure{\cdot}"] \& b'' \\
c\,{[} \arrow[r,Rightarrow,"\unsure{\cdot}"] \& c' \arrow[r,Rightarrow] \& c''{]} \& {,}
\end{tikzcd}}

\scanpage{292}
\mspage{208}
% archive sheet number: 291
où seules les implications \unsure{pointées} $b'\Rightarrow b''$ et $c\Rightarrow c'$ demandent une démonstration. Pour $b'\Rightarrow b''$, on note que l'hypothèse $C$ karoubienne implique que $C\approx M_{\unsure{\pi|\kappa|k}_0}$, \inserted{(4.8.5, $2^\circ$)} or on sait que $M_{\pi_0}$ est stable dans $M$ par \unsure{limites|plimites} inductives de cardinal $\leq\pi_0$ (et celles-ci sont représentables dans $M$ en vertu de b'), donc \deleted{\unsure{$M$ $C$ est}} $M_{\pi_0}$ \deleted{est} \deleted{$M$,} admet des $\varinjlim$ \unsure{du type envisagé|des types envisagés}, donc $C$ aussi.

$c\Rightarrow c'$ Notons que si $M=\mathrm{Ind}_{\pi_0}(C)$, alors les $\varprojlim$ représentables dans $M$ se calculent dans $C^\wedge$. Il s'ensuit que si $\pi\geq\pi_1$, alors $\mathrm{Filt}^{\pi}$ est stable par $\varprojlim$ finies ($\mathrm{Filt}^{\pi}$ \inserted{étant} donné \inserted{\unsure{alors|dans}} par la condition $a''$ de 4.7.1), d'autre part si \inserted{de plus} $\pi^{\pi_0}=\pi$, \deleted{$\mathrm{Filt}^{\pi}$} $\pi$ est un cardinal utile donc $M\approx\mathrm{Ind}_{\pi}(D)$, où $D=\mathrm{Filt}^{\pi}=M_{\pi}$. Il suffit maintenant de \unsure{voir|noter} que $\mathrm{Ind}_{\pi}(D)$ est \inserted{déjà} stable par $\varprojlim$ finies, \deleted{\unsure{pour}} \unsure{conclure} qu'il en est de $=$ de $M$, donc que $M$ est stable par limites proj. quelconques, s'il l'est par petits produits ou par petites $\varprojlim$ filtrantes, d'où $c'$. Il reste \deleted{\unsure{alors}} donc à prouver le lemme suivant \deleted{(où $\pi$ est remplacé par $\pi_0$, $D$ par $\unsure{C|D}$)}

\uline{Lemme} \unsure{4.8.10}. \deleted{\unsure{On}} Soit \deleted{$\unsure{d|D}$} $C$ une petite catégorie, $\pi_0$ un cardinal infini, $\pi_1=\mathrm{Sup}(\pi_0,\unsure{\mathrm{card\,red}|\mathrm{card}\,\text{réd}|\mathrm{card\,rel}}\,C)$ Supposons $\mathrm{Ind}_{\pi_0}(C)$ stable par $\varprojlim$ finies \inserted{(resp. \unsure{de|de $\varprojlim$} d'un certain type)} Alors \deleted{\unsure{pour}} tt $\pi\geq\pi_1$, $\mathrm{Filt}^{\pi}$ est stable par lesdites $\varprojlim$. (\unsure{Donc aussi|On a aussi} $M_\pi=\mathrm{Filt}^{\pi}$, si $\pi=\pi^{\pi_0}$). Plus\unsure{$>$|...}
